Les nombres
\[
x_1, x_2, \ldots, x_n
\]
sont écrits au tableau.

À chaque opération, on choisit deux nombres $x$ et $y$, on les efface et on
écrit à leur place le nombre
\[
x + y - xy.
\]
Après $n - 1$ opérations, il ne reste qu'un seul nombre.
\begin{enumerate}
  \item Montrer que
  \[
  1 - (x + y - xy) = (1 - x)(1 - y).
  \]
  \item Trouver une quantité multiplicative qui reste invariante au cours du
  processus.
  \item Montrer que le dernier nombre ne dépend pas de l'ordre des opérations.
  \item Exprimer ce dernier nombre en fonction de $x_1, x_2, \ldots, x_n$.
  \item Calculer le nombre final lorsque les nombres initiaux sont
  \[
  \frac{1}{2}, \frac{1}{3}, \frac{1}{4}, \ldots, \frac{1}{n+1}.
  \]
\end{enumerate}
